\chapter{Background Work}

% -----------------------------------------------------
\section{Conventional AES Architecture}
% -----------------------------------------------------

AES consists of multiple rounds that process a one hundred twenty eight bit state matrix using
a sequence of transformations. The number of rounds depends on the key size: ten for AES
one hundred twenty eight, twelve for AES one hundred ninety two, and fourteen for AES
two hundred fifty six.

\begin{figure}[!htbp]
    \centering
    \includegraphics[width=0.85\linewidth]{figs/aes-image.png}
    \caption{Conventional AES encryption and decryption round structure.}
    \label{fig:aes_conventional}
\end{figure}

The encryption process begins with an initial AddRoundKey operation followed by repeated
rounds that apply SubBytes, ShiftRows, MixColumns, and AddRoundKey. In the final round,
MixColumns is omitted. The decryption flow applies the inverse operations in reverse order.
Although the functional structure is symmetric, the conventional hardware realization requires
separate S Box and inverse S Box units, along with distinct MixColumns and InvMixColumns
modules. This leads to high area consumption and continuous switching activity across large
combinational blocks.

To summarize, conventional AES hardware commonly includes:
\begin{itemize}
    \item sixteen S Box units for SubBytes
    \item sixteen inverse S Box units for InvSubBytes
    \item four MixColumns units and four InvMixColumns units
\end{itemize}

which results in increased area and dynamic power.

% -----------------------------------------------------
\section{Proposed Architecture in the Virtex Seven Implementation}
% -----------------------------------------------------

The Virtex seven optimized AES architecture \cite{gunasekaran_aes256_virtex7} introduces two major improvements:
\begin{enumerate}
    \item a unified S Box and inverse S Box structure
    \item a combined MixColumns and AddRoundKey unit
\end{enumerate}
along with a word serial datapath that reduces the number of parallel functional units.

% --------------------------
\subsection{Unified S Box and Inverse S Box}
% --------------------------

The key observation is that both SubBytes and InvSubBytes share the same multiplicative
inverse computation, differing only in the affine transformations before and after the inverse.

\begin{figure}[!htbp]
    \centering
    \includegraphics[width=0.85\linewidth]{figs/s-box.png}
    \caption{Unified S Box and inverse S Box structure.}
    \label{fig:unified_sbox}
\end{figure}

\begin{itemize}
    \item The multiplicative inverse computation is shared between SubBytes and InvSubBytes.
    \item A mode-select multiplexer chooses between the forward and inverse affine transforms.
\end{itemize}

This unified structure eliminates duplication and reduces the number of S Box units
significantly in both parallel and serialized datapaths.

% --------------------------
\subsection{Combined MixColumns and AddRoundKey Unit}
% --------------------------

MixColumns is not required in the initial round or the final round. The optimized design merges
MixColumns and AddRoundKey into a single configurable unit that bypasses MixColumns
when the round counter indicates round zero or the last round \cite{gunasekaran_aes256_virtex7}.

\begin{figure}[!htbp]
    \centering
    \includegraphics[width=0.85\linewidth]{figs/mix-col.png}
    \caption{Combined MixColumns and AddRoundKey datapath.}
    \label{fig:mix_combined}
\end{figure}

This conditional structure:
\begin{itemize}
    \item avoids unnecessary MixColumns computation,
    \item reduces switching activity in the MixColumns module,
    \item simplifies the data path, and
    \item enables hardware reuse across rounds.
\end{itemize}

% -----------------------------------------------------
\section{Word Serial Datapath and Scheduling}
% -----------------------------------------------------

The conventional AES state matrix contains sixteen bytes arranged as a four by four grid.
Parallel architectures therefore require sixteen S Boxes and four MixColumns units. The
optimized design processes one column at a time, reducing the required functional units.

\begin{figure}[!htbp]
    \centering
    \includegraphics[width=0.45\linewidth]{figs/state_parallel.png}
    \caption{Parallel processing structure requiring sixteen S Boxes and four MixColumns units.}
    \label{fig:state_parallel}
\end{figure}

\begin{figure}[!htbp]
    \centering
    \includegraphics[width=0.45\linewidth]{figs/state_serial.png}
    \caption{Word serial processing structure requiring four S Boxes and one MixColumns unit.}
    \label{fig:state_serial}
\end{figure}

A word serial schedule processes each column in multiple cycles. An example cycle schedule
for fourteen round AES two hundred fifty six is shown below.

\begin{figure}[!htbp]
    \centering
    \includegraphics[width=0.9\linewidth]{figs/schedule_round.png}
    \caption{Example operation schedule for a word serial AES datapath.}
    \label{fig:schedule}
\end{figure}

This approach preserves correctness while reducing hardware cost and switching activity.

% -----------------------------------------------------
\section{Pipeline Architecture}
% -----------------------------------------------------

The full pipeline AES architecture implements each round as a distinct combinational stage,
enabling new plaintext to be accepted every clock cycle once the pipeline is filled.
This parallelism maximizes throughput at the expense of area: all round stages are active
simultaneously, requiring the full set of SubBytes, ShiftRows, MixColumns, and AddRoundKey
units for every pipeline stage. For AES-128 with ten rounds this results in ten independent
copies of the round logic instantiated in hardware.

In the baseline implementation, MixColumns is realized using a GF(2\textsuperscript{8})
multiplier built on EXP3 and LN3 log tables (\texttt{aes\_mcol.sv}). Each table contains
256 entries, so the combined EXP3 and LN3 lookup contributes 512 table entries to the
datapath. Because every pipeline stage contains a MixColumns unit, these tables are
replicated across all round stages, resulting in the 38,766 LUT baseline observed during
Vivado synthesis for AES-128.

The pipeline architecture is used as a high-area reference in this work. Phase~1 replaces
the log-table GF multiplier with \textit{xtime} arithmetic, eliminating the 512-entry table
from every MixColumns instance and projecting a reduction to approximately 25,000--30,000
LUTs.

% -----------------------------------------------------
\section{FSM Architecture}
% -----------------------------------------------------

The FSM based AES architecture serializes the encryption and decryption datapath across
clock cycles using an explicit state machine. A single set of round hardware is shared across
all AES rounds by cycling through states: \texttt{IDLE}, \texttt{LOAD\_KEY},
\texttt{KEY\_EXPANSION}, \texttt{ENCRYPT\_ROUND} (repeated for each round), and
\texttt{FINAL\_ROUND}. Because the same hardware block processes all rounds sequentially,
the area footprint is substantially smaller than the pipeline architecture. The AES-128
FSM baseline measures 6,794 LUTs and 412 flip-flops.

% Requires in main.tex preamble: \usepackage{tikz}
% \usetikzlibrary{automata,positioning,arrows.meta}
\begin{figure}[!htbp]
\centering
\begin{tikzpicture}[
    ->,>=stealth,
    shorten >=1pt,
    auto,
    node distance=3.2cm,
    semithick,
    every state/.style={draw, circle, minimum size=1.5cm, font=\small}
]

% States
\node[state, initial]           (IDLE)  {IDLE\\{\scriptsize state=0}};
\node[state, right of=IDLE]     (MID)   {Rounds\\{\scriptsize 1--9}};
\node[state, right of=MID]      (FINAL) {Final\\{\scriptsize state=Nr}};

% Transitions
\path (IDLE)  edge [loop above]  node {\scriptsize Enable=0}       (IDLE)
      (IDLE)  edge               node {\scriptsize Enable=1}        (MID)
      (MID)   edge [loop above]  node {\scriptsize state$<$Nr$-$1}  (MID)
      (MID)   edge               node {\scriptsize state=Nr$-$1}    (FINAL)
      (FINAL) edge [bend left=40] node {\scriptsize ready=1}        (IDLE);

\end{tikzpicture}
\caption{Encrypt FSM state diagram (AES-128, Nr=10). The CE guard
\texttt{r.state\,$\neq$\,0\,$\|$\,r.ready\,=\,1\,$\|$\,Enable\,=\,1}
freezes the packed state register when the FSM is idle, suppressing unnecessary
switching across all struct fields (Phase~2 optimization).}
\label{fig:encrypt_fsm}
\end{figure}

Three independent FSMs control the design: the encrypt FSM, the decrypt FSM, and the key
expansion FSM~\cite{mcloone_fsm_aes_2003}. Each FSM holds its computation state in a packed
register struct, enabling fine-grained clock enable gating in Phase~2.

Phase~2 of this work introduces clock enable (CE) guards on the FSM state registers. The
cipher and inverse cipher FSMs use the condition
\texttt{r.state != 0 || Enable == 1 || r.ready == 1} to gate register updates, ensuring
that the packed struct is frozen when the FSM is idle, preventing spurious switching on all
fields of the struct. The key expansion FSM uses a narrower guard
\texttt{r.state != 0 || Enable == 1} because the key expansion register has no ready bit
and does not require the ready-clear cycle present in the cipher FSMs. These CE guards
reduce dynamic switching power without altering functional behavior.

% -----------------------------------------------------
\section{Summary}
% -----------------------------------------------------

The Virtex-7 optimized AES architecture achieves substantial reduction in area and dynamic
power through structural reuse: a unified S-Box, a combined MixColumns--AddRoundKey unit, and
a word-serial datapath~\cite{gunasekaran_aes256_virtex7}.

The present work applies these principles within two distinct architectures:

\begin{itemize}
    \item \textbf{Pipeline architecture}~\cite{good_pipelined_aes_2005}: each AES round is
    instantiated as a separate combinational stage, enabling one new plaintext per clock cycle
    at the cost of replicating all round logic. Baseline AES-128 cost: 38,766 LUTs.

    \item \textbf{FSM architecture}~\cite{mcloone_fsm_aes_2003}: a single round hardware block
    is shared across all rounds under explicit state machine control. Baseline AES-128 cost:
    6,794 LUTs --- an 82\% area reduction relative to the pipeline.
\end{itemize}

Two optimization phases are applied to both architectures: Phase~1 replaces the EXP3/LN3
log-table GF multiplier with \textit{xtime} arithmetic~\cite{wolkerstorfer_xtime_2002},
targeting elimination of 512 table entries from each MixColumns instance. Phase~2 introduces
clock enable guards on the FSM state registers to reduce dynamic switching power when the
datapath is idle.